
\documentclass[a4paper,11pt]{article}

% Package imports
\usepackage{latexsym}
\usepackage{xcolor}
\usepackage{float}
\usepackage{ragged2e}
\usepackage[empty]{fullpage}
\usepackage{wrapfig}
\usepackage{lipsum}
\usepackage{tabularx}
\usepackage{titlesec}
\usepackage{geometry}
\usepackage{marvosym}
\usepackage{verbatim}
\usepackage{enumitem}
\usepackage{fancyhdr}
\usepackage{multicol}
\usepackage{graphicx}
\usepackage{cfr-lm}
\usepackage[T1]{fontenc}
\usepackage{fontawesome5}
\usepackage{soul}
\definecolor{gradhl}{RGB}{204,232,237}
\sethlcolor{gradhl}
\newcommand{\gradcourse}[1]{\hl{\textbf{#1}}}


% Color definitions
\definecolor{darkblue}{RGB}{0,0,139}

% Page layout
\setlength{\multicolsep}{0pt} 
\pagestyle{fancy}
\fancyhf{} % clear all header and footer fields
\fancyfoot{}
\renewcommand{\headrulewidth}{0pt}
\renewcommand{\footrulewidth}{0pt}
\geometry{left=1.4cm, top=0.8cm, right=1.2cm, bottom=1cm}
\setlength{\footskip}{5pt} % Addressing fancyhdr warning

% Hyperlink setup (moved after fancyhdr to address warning)
\usepackage[hidelinks]{hyperref}
\hypersetup{
    colorlinks=true,
    linkcolor=darkblue,
    filecolor=darkblue,
    urlcolor=darkblue,
}

% Custom box settings
\usepackage[most]{tcolorbox}
\tcbset{
    frame code={},
    center title,
    left=0pt,
    right=0pt,
    top=0pt,
    bottom=0pt,
    colback=gray!20,
    colframe=white,
    width=\dimexpr\textwidth\relax,
    enlarge left by=-2mm,
    boxsep=4pt,
    arc=0pt,outer arc=0pt,
}

% URL style
\urlstyle{same}

% Text alignment
\raggedright
\setlength{\tabcolsep}{0in}

% Section formatting
\titleformat{\section}{
  \scshape\raggedright\large
}{}{0em}{}[\color{black}\titlerule ]

% Custom commands
\newcommand{\resumeItem}[2]{
  \item{
    \textbf{#1}{\hspace{0.5mm}#2 }
  }
}

\newcommand{\resumePOR}[3]{
\item
    \begin{tabular*}{0.97\textwidth}[t]{l@{\extracolsep{\fill}}r}
        \textbf{#1}\hspace{0.3mm}#2 & \textit{\small{#3}} 
    \end{tabular*}
    
}

\newcommand{\resumeSubheading}[4]{
\item
    \begin{tabular*}{0.98\textwidth}[t]{l@{\extracolsep{\fill}}r}
        \textbf{#1} & \textit{\footnotesize{#4}} \\
        \textit{\footnotesize{#3}} &  \footnotesize{#2}\\
    \end{tabular*}
    
}

\newcommand{\resumeProject}[4]{
\item
    \begin{tabular*}{0.98\textwidth}[t]{l@{\extracolsep{\fill}}r}
        \textbf{#1} & \textit{\footnotesize{#3}} \\
        \footnotesize{\textit{#2}} & \footnotesize{#4}
    \end{tabular*}
    
}

\newcommand{\resumeSubItem}[2]{\resumeItem{#1}{#2}}

\renewcommand{\labelitemi}{$\vcenter{\hbox{\tiny$\bullet$}}$}
\renewcommand{\labelitemii}{$\vcenter{\hbox{\tiny$\circ$}}$}

\newcommand{\resumeSubHeadingListStart}{\begin{itemize}[leftmargin=*,labelsep=1mm]}
\newcommand{\resumeHeadingSkillStart}{\begin{itemize}[leftmargin=*,itemsep=1.7mm, rightmargin=2ex]}
\newcommand{\resumeItemListStart}{\begin{itemize}[leftmargin=*,labelsep=1mm,itemsep=0.5mm]}

\newcommand{\resumeSubHeadingListEnd}{\end{itemize}}
\newcommand{\resumeHeadingSkillEnd}{\end{itemize}}
\newcommand{\resumeItemListEnd}{\end{itemize}}
\newcommand{\cvsection}[1]{%

\begin{tcolorbox}
    \textbf{\large #1}
\end{tcolorbox}
    
}

\newcolumntype{L}{>{\raggedright\arraybackslash}X}%
\newcolumntype{R}{>{\raggedleft\arraybackslash}X}%
\newcolumntype{C}{>{\centering\arraybackslash}X}%

% Commands for icon sizing and positioning
\newcommand{\socialicon}[1]{\raisebox{-0.05em}{\resizebox{!}{1em}{#1}}}
\newcommand{\ieeeicon}[1]{\raisebox{-0.3em}{\resizebox{!}{1.3em}{#1}}}

% Font options
\newcommand{\headerfonti}{\fontfamily{phv}\selectfont} % Helvetica-like (similar to Arial/Calibri)
\newcommand{\headerfontii}{\fontfamily{ptm}\selectfont} % Times-like (similar to Times New Roman)
\newcommand{\headerfontiii}{\fontfamily{ppl}\selectfont} % Palatino (elegant serif)
\newcommand{\headerfontiv}{\fontfamily{pbk}\selectfont} % Bookman (readable serif)
\newcommand{\headerfontv}{\fontfamily{pag}\selectfont} % Avant Garde-like (similar to Trebuchet MS)
\newcommand{\headerfontvi}{\fontfamily{cmss}\selectfont} % Computer Modern Sans Serif
\newcommand{\headerfontvii}{\fontfamily{qhv}\selectfont} % Quasi-Helvetica (another Arial/Calibri alternative)
\newcommand{\headerfontviii}{\fontfamily{qpl}\selectfont} % Quasi-Palatino (another elegant serif option)
\newcommand{\headerfontix}{\fontfamily{qtm}\selectfont} % Quasi-Times (another Times New Roman alternative)
\newcommand{\headerfontx}{\fontfamily{bch}\selectfont} % Charter (clean serif font)

\begin{document}
\headerfontiii

% Header
\begin{center}
    {\Huge\textbf{David Xia}}
\end{center}

\begin{center}
    \href{mailto:davidx3@illinois.edu}{davidx3@illinois.edu} | \href{https://davidxia3.github.io/index.html}{Personal Website} | 
    \socialicon{\faLinkedin} \href{https://www.linkedin.com/in/dx3/}{LinkedIn - dx3} | 
    \socialicon{\faGithub} \href{https://github.com/davidxia3}{GitHub - davidxia3}
\end{center}

\section{\textbf{Research Interests}}
Reinforcement learning, foundation models, LLM agents, RL theory

\section{\textbf{Education}}
\resumeSubHeadingListStart
\resumeSubheading
{University of Illinois Urbana-Champaign}{Champaign, Illinois}
{\parbox[t]{6in}{Bachelor of Science in Computer Science \\ Bachelor of Science in Liberal Arts \& Sciences Major in Mathematics, Data Optimization Concentration\\
Bachelor of Science in Liberal Arts \& Sciences Major in Statistics}}
{\hspace{-3ex}Aug 2023 - May 2027}
\resumeItemListStart
\item GPA: 3.99/4.00
\resumeItemListEnd
\resumeSubHeadingListEnd

\section{\textbf{Papers} \hfill \textcolor{darkblue}{\scriptsize REV=Under Review, ACC=Accepted}}

\begin{enumerate}[leftmargin=*, labelsep=0.5em, align=left, widest={[\textbf{S.1}]}, itemindent=0em, label={\textbf{[\arabic*]}}]
    \item[\textbf{2026}] \textbf{[REV]} \textbf{Benchmarking agentic optimizers}\\
        Co-first author. Title and author list withheld during double-blind review.\\
        \textit{Under review at ICLR 2027}
    \item[\textbf{2026}] \textbf{[ACC]} \textbf{RSK linear operators and the Vershik-Kerov-Logan-Shepp curve} [\href{https://www.combinatorics.org/ojs/index.php/eljc/article/view/v33i3p17}{Journal}] [\href{https://arxiv.org/abs/2504.05510}{arXiv}]\\ 
        Duy Phan, \textbf{David Xia}\\
        \textit{Electronic Journal of Combinatorics, 33(3), \#P17}
    \item[\textbf{2025}] \textbf{[ACC]} \textbf{EveGuard: Defeating Vibration-based Side-Channel Eavesdropping with Audio Adversarial Perturbations} [\href{https://ieeexplore.ieee.org/document/11023366/}{IEEE}] [\href{https://arxiv.org/abs/2411.10034}{arXiv}]\\ 
        Jung-Woo Chang, Ke Sun, \textbf{David Xia}, Xinyu Zhang, Farinaz Koushanfar \\
        \textit{IEEE Symposium on Security and Privacy, 2025}
\end{enumerate}

\section{\textbf{Workshop Papers, Manuscripts \& Notes} \hfill \textcolor{darkblue}{\scriptsize NOTE=Technical Note, SUB=Submitted, ACC=Accepted}}

\begin{enumerate}[leftmargin=*, labelsep=0.5em, align=left, widest={[\textbf{S.1}]}, itemindent=0em, label={\textbf{[\arabic*]}}]
    \item[\textbf{2026}] \textbf{[NOTE]} \textbf{Frozen Policy Iteration for MDPs with Stochastic Transitions under Linear $Q^\pi$ Realizability} [\href{https://davidxia3.github.io/docs/stochastic_fpi.pdf}{PDF}] \\
        \textbf{David Xia}, Ruizhong Qiu, Hanghang Tong\\
        \textit{Technical note}
    \item[\textbf{2026}] \textbf{[SUB]} \textbf{Predictability in Major League Sports: Betting Odds versus Mathematical Models} [\href{https://davidxia3.github.io/docs/predictability_paper.pdf}{PDF}] \\
        \textbf{David Xia}, AJ Hildebrand\\
        \textit{Submitted to the Mathematics and Sports (MAS) Journal}
    \item[\textbf{2026}] \textbf{[ACC]} \textbf{Revitalizing Local Democracy: A Human-Centered Audit of LLMs in City Council Journalism} [\href{https://davidxia3.github.io/docs/heal_paper.pdf}{PDF}]\\
        \textbf{David Xia}, Chris Maury\\
        \textit{Human-centered Evaluation and Auditing of Language Models (HEAL) Workshop at CHI 2026}
\end{enumerate}

\section{\textbf{Research Experience}}
  \resumeSubHeadingListStart

    \resumeSubheading
      {{\href{https://zhaoranwang.github.io/}{Northwestern University}}\textnormal{, Research Intern}}{}
      {}{May 2026 - present}
      \resumeItemListStart
        \vspace{-2em}
        \item Advised by Professor \href{https://zhaoranwang.github.io/}{Zhaoran Wang}.
        \item Research on LLM agents as black-box optimizers: how agents propose candidate solutions, interpret feedback, and revise their search strategy across iterations.
        \item Co-first-authored paper under double-blind review at ICLR 2027. Further details withheld during the review period.
      \resumeItemListEnd

    \resumeSubheading
      {{\href{https://www.idea-isail.com/index.html}{iDEA-iSAIL Lab}}\textnormal{, Research Intern}}{}
      {}{Feb 2026 - present}
      \resumeItemListStart
        \vspace{-2em}
        \item Advised by Professor \href{http://tonghanghang.org/}{Hanghang Tong}.
        \item Conducted research in theoretical reinforcement learning, studying sample and computation efficient algorithms for online learning in finite-horizon Markov Decision Processes under function approximation and stochastic dynamics.
        \item Used concentrability to develop the first sample and computation efficient algorithm for online reinforcement learning in stochastic-transition MDPs under the linear $Q^\pi$ realizability assumption, resolving an open problem left by prior work restricted to deterministic dynamics.
        \item Established a PAC sample complexity guarantee of $\tilde{O}(d^2H^7C^*/\varepsilon^3)$ by combining techniques from both offline and online reinforcement learning theory, dynamic programming, concentration analysis, and linear function approximation.
        \item \href{https://davidxia3.github.io/docs/stochastic_fpi.pdf}{Technical note.}
      \resumeItemListEnd 

    \resumeSubheading
      {{\href{https://ayong.web.illinois.edu/ICLUE.html}{Illinois Combinatorics Lab for Undergraduate Experience (ICLUE)}}\textnormal{, Research Intern}}{}
      {}{Aug 2024 - May 2026}
      \resumeItemListStart
        \vspace{-2em}
        \item Advised by Professor \href{https://ayong.web.illinois.edu/}{Alexander Yong}.
        \item Proved a new asymptotic result that the probability that the Schensted insertion algorithm for a uniformly random permutation in $S_n$ exhibits a special ``bumping'' interaction converges to 1 as $n$ approaches infinity.
        \item Derived the result using rigorous probability bounds on the RSK linear operator by leveraging the Vershik-Kerov-Logan-Shepp limit shape curve as the central analytic tool.
        \item Developed proofs requiring techniques spanning real analysis (uniform convergence, limit arguments), probability theory (concentration inequalities, convergence in probability), and asymptotic analysis (bounding combinatorial quantities in the large-$n$ regime).
        \item Paper published in the \href{https://www.combinatorics.org/ojs/index.php/eljc/article/view/v33i3p17}{Electronic Journal of Combinatorics}, 33(3), \#P17.
      \resumeItemListEnd 

    \resumeSubheading
      {{\href{https://iml.math.illinois.edu/}{Illinois Mathematics Lab}}\textnormal{, Research Intern}}{}
      {}{Jan 2024 - May 2026}
      \resumeItemListStart
        \vspace{-2em}
        \item Advised by Professor AJ Hildebrand.
        \item Researched the data science perspective of sports matches, analyzing the predictability of outcomes by comparing prediction accuracies across polls, the betting market, Elo ratings, and mathematical models like Bradley-Terry.
        \item Presented findings at Rose-Hulman Undergraduate Math Conference 2024, UIUC Undergraduate Research Symposium 2024, and Joint Mathematics Meeting 2025.
        \item Paper in submission to the \href{https://janeway.uncpress.org/ms/}{Mathematics and Sports (MAS) Journal}.
      \resumeItemListEnd 

    \resumeSubheading
      {{\href{https://hcii.cmu.edu/}{Carnegie Mellon University, HCII}}\textnormal{, Research Intern}}{}
      {}{May 2025 - Feb 2026}
      \resumeItemListStart
        \vspace{-2em}
        \item Advised by Professor \href{https://www.cs.cmu.edu/~jbigham/}{Jeff Bigham}.
        \item Designed and evaluated an automated journalism pipeline using SOTA LLMs to replicate end-to-end editorial workflows, from transcript segmentation to headline generation and topic prioritization.
        \item Orchestrated large-scale crowd-sourced study and demonstrated that LLM-based headline quality and topic prioritization can meet and exceed professional standards.
        \item Paper in CHI 2026 \href{https://heal-workshop.github.io/}{Human-centered Evaluation and Auditing of Language Models (HEAL) workshop}.
      \resumeItemListEnd

    \resumeSubheading
      {{\href{https://blog.nus.edu.sg/seeder/}{National University of Singapore, SERIUS REU}}\textnormal{, Research Intern}}{}
      {}{May 2024 - Aug 2024}
      \resumeItemListStart
        \vspace{-2em}
        \item Advised by Professor Kelvin Fong Xuanyao.
        \item Designed the GARDEN (Generalized Anomaly Recognition and Detection for Enhanced Nurturing) pipeline for early disease detection in farm plants, combining segmentation, generative, and classification models end-to-end.
        \item Built a leaf segmentation model to isolate regions of interest from complex backgrounds, enabling downstream anomaly detection to focus on plant tissue only.
        \item Implemented a conditional GAN (cGAN) to synthesize realistic diseased and healthy leaf images, augmenting training data and improving classifier generalization.
        \item Achieved robust disease classification by ensembling deep learning and traditional ML models, exploiting complementary strengths to improve detection reliability on real farm images.
      \resumeItemListEnd

    \resumeSubheading
      {{\href{https://xyzhang.ucsd.edu/index.html}{University of California San Diego}}\textnormal{, Research Intern}}{}
      {}{Jun 2023 - Dec 2024}
      \resumeItemListStart
        \vspace{-2em}
        \item Advised by Professor \href{https://xyzhang.ucsd.edu/index.html}{Xinyu Zhang}.
        \item Developed EveGuard, a software-based defense framework to protect voice privacy from vibrometry-based side channels, using a perturbation generator model (PGM) to suppress sensor-based eavesdropping while preserving high audio quality.
        \item Implemented Eve-GAN, a novel domain-translation task for inferring eavesdropped signals, enabling end-to-end training of the PGM and using few-shot learning to reduce data-collection overhead.
        \item Achieved a protection rate of over 97\% against audio classifiers, hindering eavesdropped reconstruction.
        \item Paper published in the \href{https://ieeexplore.ieee.org/document/11023366/}{IEEE Symposium on Security and Privacy 2025}.
      \resumeItemListEnd

  \resumeSubHeadingListEnd

\section{\textbf{Relevant Coursework}}

\begin{itemize}[leftmargin=0.15in, label={}, nosep]
\item{
    {\footnotesize Highlighted courses are graduate-level.} \\[0.6em]
    \textbf{In progress (Fall 2026):} \gradcourse{Statistical Reinforcement Learning}, \gradcourse{Randomized Algorithms} \\
    \textbf{CS:} Machine Learning, Deep Learning, Natural Language Processing, Algorithms, Numerical Analysis, Numerical Methods \\
    \textbf{Math:} Honors Real Analysis, \gradcourse{Algebraic Combinatorics}, Abstract Linear Algebra, Combinatorics, Graph Theory, Abstract Algebra, Complex Variables \\
    \textbf{Stats:} Bayesian Analysis
  }
\end{itemize}

\section{\textbf{Grants and Awards}}

\begin{enumerate}[leftmargin=*, labelsep=0.5em, align=left, widest={[\textbf{S.1}]}, itemindent=0em, label={\textbf{[\arabic*]}]}]
\item[\textbf{2025}] CMU HCII REU: NSF 2349558
\item[\textbf{2025}] ICLUE RTG: NSF 1937241
\end{enumerate}

\begin{center}
    \href{mailto:davidx3@illinois.edu}{davidx3@illinois.edu} | \href{https://davidxia3.github.io/index.html}{Personal Website} | 
    \socialicon{\faLinkedin} \href{https://www.linkedin.com/in/dx3/}{LinkedIn - dx3} | 
    \socialicon{\faGithub} \href{https://github.com/davidxia3}{GitHub - davidxia3}
\end{center}

\end{document}
